\section{Conclusion and Remaining Work}\label{sec:conclusion}
We asked whether the advantage scMINER reports for mutual information clustering survives when
the method must choose its own clustering without reference labels. Under one published
label-free rule, applied identically to MICA and to a principal component and Leiden pipeline
on ten datasets with repeats, part of it does. On the nine small datasets MICA leads under three
of four rules by small margins, and when its saved scores are restricted after the fact to the
cluster counts the baseline produced it leads under both readings. The size of that lead depends
on the rule's
tie-break, threshold and random draw, and most of all on which counts each method's grid
offers, which a cap of 20 clusters does not equalise. The same confound accounts for the
difference we first saw between MICA's two pipelines, and on PBMC14K, the one large benchmark dataset,
the ranking turns on the threshold. These findings are about one rule, and we do not extend
them to mutual information clustering under other rules.

The faster neighbour graph is a faster graph rather than a better one. It gives the same
recall and clustering as MICA's on PBMC14K at a fraction of the build time, and on the Allen
cells higher recall in less than half the build time, without a better clustering. On both it needs
three to four times the memory.

\subsection*{Plan through November 4}
The main experiments are due on November 4 and the complete draft on November 18. Three items
remain before November 4, and each answers a limitation above rather than adding scope.

\textbf{Matched candidate grids.} Extend the baseline's resolution grid until it reaches 2 to
20 clusters on every dataset, and MICA's graph-mode grid until it reaches 2, then rerun the
stability comparison and the mode test. The common-range result above comes from restricting
saved scores after the fact, and it should rest on runs designed for it.

\textbf{A transfer comparison, reported apart from the label-free rules.} Each dataset takes
the setting that scored best on the other nine. Because that choice uses the other datasets'
labels it is not label-free, and it answers a different question: how far a setting tuned on
labelled benchmarks carries to a new one. It also tests whether MICA's shipped default of
four clusters behaves like such a tuned setting.

\textbf{The final comparison table} across the ten datasets, every rule with its repeats and
costs, with no win-count criterion. The question is how the ranking moves with the rule, not
whether MICA reaches a fixed number of wins.

\textbf{Deferred until those are complete:} cluster annotation with a language model, the hub
correction for the mutual information graph, and GPU work, as the feedback advised.

\textbf{The full atlas.} Our 100,000-cell pilot of MICA as shipped took 7.1 hours and peaked
at 140 GB, which scales to about 6 to 8 days and 3.2 TB for all 2.3 million cells, against the
754 GB on our machine. A faster graph does not close that gap. Our shortlist build alone
peaked at 80 GB on the same cells, about 1.8 TB in proportion, so the whole atlas would need a
sketch, clustering a sample and assigning the rest. We will not change the course-scale plan
for the atlas without first discussing it with the instructor.

\section*{Completion plan}
\textbf{Utshab} reruns the stability comparison on matched grids, runs the transfer comparison
and builds the final table, drafts the methods, design and results sections, and raises the
full-atlas question with the instructor. All computation runs on one workstation, which only Utshab has access to.

\textbf{Laraib} completes the dataset cards for PBMC14K and the Allen subset, finishes the
related work comparison table and revises the introduction and related work from it, and
checks independently that the interpretation in the discussion follows from the exported
results.

\textbf{Rhea} regenerates the result tables as new runs land, checks every number in the text
against the exported results without relying on the author's reading of them, keeps the
schedule and risk list current, and assembles the slides. Rhea raised six anomalies against
the Assignment 2 tables, all of which are resolved in Appendix~\ref{app:anomalies}.
